\chapter{Evaluierung und Ergebnisse}
\abschnittsautor{\alleautoren}

\section{Soll-Ist-Abgleich}

Stellt die Anforderungen aus der Zielsetzung dem tatsächlichen Ergebnis
gegenüber. Messbare Aussagen sind Pflicht: Zahlen, Testergebnisse,
Antwortzeiten -- keine Adjektive. Tab.~\ref{tab:sollist} zeigt den Abgleich.

Auf jede Tabelle und jede Abbildung muss im Text Bezug genommen werden. Eine
Tabelle, die nirgends erwähnt wird, gilt als nicht erklärt.

\begin{table}[htbp]
  \caption{Soll-Ist-Abgleich der Anforderungen}
  \label{tab:sollist}
  \centering
  \begin{tabular}{@{}p{0.45\linewidth}p{0.2\linewidth}p{0.25\linewidth}@{}}
    \toprule
    \bfseries Anforderung & \bfseries Soll & \bfseries Ist \\
    \midrule
    Antwortzeit der Produktliste & < 500\,ms & 380\,ms \\
    Abweisung fremder Konten     & 100\,\%   & 100\,\% \\
    \bottomrule
  \end{tabular}
\end{table}

\section{Fazit und Ausblick}

Was ist fertig, was fehlt, und was wären die nächsten sinnvollen Schritte?
